\clearpage\corrigesChapitre{5}{Électrons et stabilité}
\corrigesSouples
\begin{blocLivre}{cor-5-1}\begin{corrige}{1}{Additionner les exposants}
\schema[corrige]{tableau-periodique}
\begin{questions}
\item $2+2+5=9$ électrons, donc $Z=9$.
\item C’est le fluor F.
\end{questions}
\end{corrige}\end{blocLivre}
\begin{blocLivre}{cor-5-2}\begin{corrige}{2}{Électrons de valence}
\begin{questions}
\item Période 3, deux électrons de valence dans la couche externe.
\item La perte de deux électrons conduit à $\ce{Mg^{2+}}$, de configuration externe identique à celle du néon.
\end{questions}
\end{corrige}\end{blocLivre}
\begin{blocLivre}{cor-5-3}\begin{corrige}{3}{Nom et formule}
\begin{questions}
\item $\ce{Na+}$, $\ce{Cl-}$, $\ce{Ca^{2+}}$, $\ce{F-}$.
\item Sodium et calcium : cations ; chlorure et fluorure : anions.
\end{questions}
\end{corrige}\end{blocLivre}
\begin{blocLivre}{cor-5-4}\begin{corrige}{4}{Interpréter l’eau}
\schema[corrige]{lewis}
\begin{questions}
\item L’oxygène possède 2 doublets liants et 2 non liants, soit 8 électrons autour de lui.
\item Chaque H partage un doublet liant, donc est entouré de 2 électrons.
\end{questions}
\end{corrige}\end{blocLivre}
\begin{blocLivre}{cor-5-5}\begin{corrige}{5}{Écrire une configuration}
\begin{questions}
\item $1s^2\,2s^2\,2p^6\,3s^2\,3p^2$ ; les exposants somment à 14.
\item Période 3 ; 4 électrons de valence.
\end{questions}
\end{corrige}\end{blocLivre}
\begin{blocLivre}{cor-5-6}\begin{corrige}{6}{Deux membres d’une famille}
\begin{questions}
\item Leur couche externe contient six électrons, sous la forme $s^2p^4$.
\item Cette structure de valence commune explique leur appartenance à une même famille et certaines propriétés analogues.
\end{questions}
\end{corrige}\end{blocLivre}
\begin{blocLivre}{cor-5-7}\begin{corrige}{7}{Un octet dans le dioxyde de carbone}
\schema[corrige]{lewis}
\begin{questions}
\item Deux liaisons doubles correspondent à quatre doublets liants autour du carbone.
\item Carbone : $4\times2=8$ électrons. Chaque oxygène a 2 doublets liants et 2 non liants, donc 8 électrons.
\end{questions}
\end{corrige}\end{blocLivre}
\begin{blocLivre}{cor-5-8}\begin{corrige}{8}{Comparer des liaisons}
\begin{questions}
\item La seconde : son énergie de rupture est plus grande.
\item Non, rompre une liaison nécessite un apport d’énergie. C’est la formation d’une liaison qui en libère.
\end{questions}
\end{corrige}\end{blocLivre}
\begin{blocLivre}{cor-5-9}\begin{corrige}{9}{Un gaz noble particulier}
\begin{questions}
\item Leur couche externe est complète dans le modèle étudié.
\item Non : la couche 1 ne peut contenir que deux électrons. Le duet suffit pour l’hélium.
\end{questions}
\end{corrige}\end{blocLivre}
\begin{blocLivre}{cor-5-10}\begin{corrige}{10}{Corriger un remplissage}
\begin{questions}
\item La sous-couche 2p ne peut contenir que six électrons.
\item $1s^2\,2s^2\,2p^6\,3s^1$ ; l’ion est $\ce{Na+}$.
\end{questions}
\end{corrige}\end{blocLivre}
\begin{blocLivre}{cor-5-11}\begin{corrige}{11}{La molécule d’ammoniac}
\schema[corrige]{lewis}
\begin{questions}
\item 3 doublets liants et 1 non liant : 8 électrons entourent l’azote.
\item Il y a 4 doublets au total, soit 8 électrons de valence représentés. Les doublets partagés ne se comptent qu’une fois dans le total moléculaire.
\end{questions}
\end{corrige}\end{blocLivre}
\begin{blocLivre}{cor-5-12}\begin{corrige}{12}{Atome ou ion ?}
\demarche{Comparer séparément le nombre de protons et celui des électrons.}
\begin{questions}
\item 10 protons : atome de néon ; 12 protons et 10 électrons : ion $\ce{Mg^{2+}}$.
\item Les deux configurations sont $1s^2\,2s^2\,2p^6$. Même cortège ne signifie pas même élément.
\end{questions}
\retiens{Deux entités peuvent posséder le même cortège électronique sans être le même élément.}
\end{corrige}\end{blocLivre}

